\documentclass[11pt]{article}

\usepackage[T1]{fontenc}
\usepackage[utf8]{inputenc}
\usepackage[spanish,es-nodecimaldot]{babel}
\usepackage{amsmath,amssymb,mathtools}
\usepackage{enumitem}
\usepackage[a4paper,margin=2.5cm]{geometry}
\usepackage[hidelinks]{hyperref}

\newcommand{\Asa}{A_{\mathrm{sa}}}
\newcommand{\U}{\mathcal{U}}
\newcommand{\Pj}{\mathcal{P}}
\newcommand{\T}{\mathbb{T}}
\newcommand{\GL}{\mathrm{GL}}
\newcommand{\id}{\mathrm{id}}

\newenvironment{teorema}[1]
  {\par\medskip\noindent\textbf{Teorema #1.}\ \itshape}
  {\par\medskip}
\newenvironment{proposicion}[1]
  {\par\medskip\noindent\textbf{Proposición #1.}\ \itshape}
  {\par\medskip}
\newenvironment{corolario}[1]
  {\par\medskip\noindent\textbf{Corolario #1.}\ \itshape}
  {\par\medskip}
\newenvironment{lema}[1]
  {\par\medskip\noindent\textbf{Lema #1.}\ \itshape}
  {\par\medskip}
\newenvironment{observacion}[1]
  {\par\medskip\noindent\textbf{Observación #1.}\ }
  {\par\medskip}

\title{Enunciados de la Parte 2: $C^*$-álgebras}
\author{}
\date{}

\begin{document}
\maketitle

\section*{Índice}
\begin{itemize}[leftmargin=*]
  \item \textbf{Observaciones:} 2.2, 2.9, 2.24, 2.25.
  \item \textbf{Lemas:} 2.10, 2.26.
  \item \textbf{Proposiciones:} 2.11, 2.12, 2.13, 2.16, 2.17, 2.18, 2.28, 2.29, 2.30.
  \item \textbf{Corolarios:} 2.7, 2.14, 2.15, 2.21.
  \item \textbf{Teoremas:} 2.6, 2.8, 2.19, 2.20, 2.22, 2.27.
\end{itemize}

\section*{Enunciados}

% Fuente: Parte2.pdf, p. 1
\begin{observacion}{2.2}
Si $A$ tiene unidad, entonces
\[
  1^*=1
  \qquad\text{y}\qquad
  \lVert 1\rVert=1.
\]
\end{observacion}

% Fuente: Parte2.pdf, p. 1
\begin{teorema}{2.6}
Sea $A$ una $C^*$-álgebra con unidad y $a\in A$ un elemento normal. Entonces
\[
  \lVert a\rVert=r(a).
\]
\end{teorema}

% Fuente: Parte2.pdf, p. 2
\begin{corolario}{2.7}
La norma en una $C^*$-álgebra es única y está dada por
\[
  \lVert a\rVert=r(a^*a)^{1/2}.
\]
\end{corolario}

% Fuente: Parte2.pdf, p. 2
\begin{teorema}{2.8}
Sean $A$ y $B$ dos $C^*$-álgebras con unidad y $\Phi:A\to B$ un $*$-homomorfismo unital. Entonces $\Phi$ es acotado y
\[
  \lVert\Phi\rVert=1.
\]
\end{teorema}

% Fuente: Parte2.pdf, p. 2
\begin{observacion}{2.9}
Sea $a\in A$. Entonces existen elementos únicos $x,y\in\Asa$ tales que
\[
  a=x+iy.
\]
Precisamente,
\[
  x=\frac{a+a^*}{2},
  \qquad
  y=\frac{a-a^*}{2i}.
\]
Llamamos a $x$ la \emph{parte real} y a $y$ la \emph{parte imaginaria} de $a$ (a veces $\operatorname{re}(a)=x$, $\operatorname{im}(a)=y$).
\end{observacion}

% Fuente: Parte2.pdf, p. 2
\begin{lema}{2.10}
Sea $A$ una $C^*$-álgebra con unidad y $x\in\Asa$. Entonces
\[
  \lVert x+it1\rVert^2\leq \lVert x\rVert^2+t^2,
  \qquad t\in\mathbb{R}.
\]
\end{lema}

% Fuente: Parte2.pdf, pp. 2--3
\begin{proposicion}{2.11}
Sea $z\in\mathbb{C}$. Supongamos que existe $C\geq0$ tal que
\[
  |z+it|^2\leq C+t^2
  \qquad \forall t\in\mathbb{R}.
\]
Entonces $z\in\mathbb{R}$.
\end{proposicion}

% Fuente: Parte2.pdf, p. 3
\begin{proposicion}{2.12}
Sea $A$ una $C^*$-álgebra con unidad y $\varphi:A\to\mathbb{C}$ un funcional lineal tal que
\[
  \lVert\varphi\rVert=\varphi(1).
\]
Entonces
\[
  \varphi(a)\in\mathbb{R}
  \qquad\text{para todo } a\in\Asa.
\]
\end{proposicion}

% Fuente: Parte2.pdf, p. 3
\begin{proposicion}{2.13}
Sea $A$ una $C^*$-álgebra con unidad y $a\in\Asa$. Entonces
\[
  \sigma(a)\subseteq\mathbb{R}.
\]
\end{proposicion}

% Fuente: Parte2.pdf, p. 3
\begin{corolario}{2.14}
Sea $A$ una $C^*$-álgebra con unidad y $\varphi:A\to\mathbb{C}$ un funcional lineal tal que
\[
  \lVert\varphi\rVert=\varphi(1).
\]
Entonces
\[
  \varphi(a^*)=\overline{\varphi(a)}
  \qquad\text{para todo }a\in A.
\]
En particular, los caracteres lo cumplen.
\end{corolario}

% Fuente: Parte2.pdf, pp. 3--4
\begin{corolario}{2.15}
Sea $A$ una $C^*$-álgebra conmutativa con unidad. Entonces la transformada de Gelfand
\[
  \Lambda:A\longrightarrow C(\Omega(A)),
  \qquad
  a\longmapsto \widehat a,
\]
es un $*$-homomorfismo unital.
\end{corolario}

% Fuente: Parte2.pdf, p. 4
\begin{proposicion}{2.16}
Sea $A$ una $C^*$-álgebra con unidad y $a\in A$. Entonces
\[
  a\in\GL(A)
  \quad\Longleftrightarrow\quad
  a^*a,\;aa^*\in\GL(A).
\]
\end{proposicion}

% Fuente: Parte2.pdf, p. 4
\begin{proposicion}{2.17}
Sea $A$ una $C^*$-álgebra con unidad $1_A$ y $B\subseteq A$ una $C^*$-subálgebra que contiene a la unidad $1_A$. Sea $b\in B$. Entonces
\[
  b\in\GL(B)
  \quad\Longleftrightarrow\quad
  b\in\GL(A),
\]
y en particular
\[
  \sigma_B(b)=\sigma_A(b).
\]
\end{proposicion}

% Fuente: Parte2.pdf, pp. 4--5
\begin{proposicion}{2.18}
Sea $A$ una $C^*$-álgebra con unidad y $a\in A$ normal. Entonces la función
\[
  \Omega\bigl(C^*(a,1)\bigr)\longrightarrow\sigma(a),
  \qquad
  \varphi\longmapsto\varphi(a),
\]
es un homeomorfismo.
\end{proposicion}

% Fuente: Parte2.pdf, p. 5
\begin{teorema}{2.19 (Stone--Weierstrass)}
Sea $X$ un espacio compacto y Hausdorff, y sea $A\subseteq C(X)$ una $*$-subálgebra con unidad. Supongamos que $A$ separa puntos, i.e. si $x\neq y$ entonces existe $f\in A$ con $f(x)\neq f(y)$. Entonces $A$ es densa en $C(X)$ con la norma $\lVert\cdot\rVert_\infty$.
\end{teorema}

% Fuente: Parte2.pdf, p. 5
\begin{teorema}{2.20 (Gelfand--Naimark)}
Sea $A$ una $C^*$-álgebra conmutativa con unidad. Entonces la transformada de Gelfand
\[
  \Lambda:A\longrightarrow C(\Omega(A)),
  \qquad
  \widehat a(\varphi)=\varphi(a),
\]
es un $*$-isomorfismo isométrico.
\end{teorema}

% Fuente: Parte2.pdf, p. 6
\begin{corolario}{2.21}
Toda $C^*$-álgebra conmutativa con unidad es isomorfa a $C(X)$ para algún espacio compacto y Hausdorff $X$. Además, $X$ es único salvo homeomorfismo.
\end{corolario}

% Fuente: Parte2.pdf, p. 6
\begin{teorema}{2.22 (Cálculo Funcional)}
Sea $A$ una $C^*$-álgebra con unidad y $a\in A$ normal. Entonces existe un único $*$-homomorfismo isométrico
\[
  C(\sigma(a))\longrightarrow A,
  \qquad
  f\longmapsto f(a),
\]
tal que
\[
  \id_{\sigma(a)}\longmapsto a
  \qquad\text{y}\qquad
  1_{\sigma(a)}\longmapsto 1_A.
\]
\end{teorema}

% Fuente: Parte2.pdf, p. 7
\begin{observacion}{2.24}
Por ejemplo,
\[
  (\overline z z)(a)=a^*a,
  \qquad
  (z^2-\overline z^{\,3}+1)(a)=a^2-(a^*)^3+1,
  \qquad\text{etc.}
\]
\end{observacion}

% Fuente: Parte2.pdf, p. 7
\begin{observacion}{2.25}
Supongamos que $X$ es compacto y $A:=C(X)$. Fijemos $g\in C(X)$, cuyo espectro es
\[
  \sigma(g)=g(X).
\]
El cálculo funcional está dado por
\[
  C(g(X))\longrightarrow A,
  \qquad
  f\longmapsto f(g)=f\circ g.
\]
Se sigue que lo anterior define un $*$-homomorfismo unital que cumple
\[
  \id_{g(X)}\longmapsto g.
\]
\end{observacion}

% Fuente: Parte2.pdf, p. 7
\begin{lema}{2.26}
Sean $A$ y $B$ dos $C^*$-álgebras con unidad y $\Phi:A\to B$ un $*$-homomorfismo unital. Sea $a\in A$ un elemento normal y $f:\sigma(a)\to\mathbb{C}$ una función continua. Entonces
\[
  \Phi\bigl(f(a)\bigr)=f\bigl(\Phi(a)\bigr).
\]
\end{lema}

% Fuente: Parte2.pdf, p. 7
\begin{teorema}{2.27}
Sean $A$ y $B$ dos $C^*$-álgebras con unidad y $\Phi:A\to B$ un $*$-homomorfismo inyectivo unital. Entonces $\Phi$ es isométrico, i.e.
\[
  \lVert\Phi(a)\rVert=\lVert a\rVert
  \qquad\text{para todo }a\in A.
\]
\end{teorema}

% Fuente: Parte2.pdf, p. 8
\begin{proposicion}{2.28}
Sea $A$ una $C^*$-álgebra con unidad y $a\in A$ un elemento normal. Entonces
\begin{enumerate}[label=(\roman*),leftmargin=2.2em]
  \item $a=a^* \Longleftrightarrow \sigma(a)\subseteq\mathbb{R}$;
  \item $a\in\U(A) \Longleftrightarrow \sigma(a)\subseteq\T$;
  \item $a\in\Pj(A) \Longleftrightarrow \sigma(a)\subseteq\{0,1\}$.
\end{enumerate}
\end{proposicion}

% Fuente: Parte2.pdf, p. 8
\begin{proposicion}{2.29}
Sea $A$ una $C^*$-álgebra con unidad. Entonces cada elemento de $A$ es combinación lineal de a lo más cuatro unitarios.
\end{proposicion}

% Fuente: Parte2.pdf, p. 8
\begin{proposicion}{2.30}
Sea $A$ una $C^*$-álgebra con unidad y $a\in A$ normal con
\[
  \sigma(a)\subseteq\mathbb{R}^+=[0,\infty).
\]
Entonces existe $y=y^*\in A$ con $\sigma(y)\subseteq\mathbb{R}^+$ tal que
\[
  y^2=a.
\]
\end{proposicion}

\end{document}
